
\section{Técnicas, instrumentos y plan de análisis}

\subsection{Técnica principal: encuesta}

Se propone una encuesta estructurada y autoadministrada, preferiblemente digital y accesible desde computador o teléfono. El instrumento incluirá una explicación breve del estudio, el consentimiento informado, los criterios de inclusión y preguntas sobre las herramientas digitales utilizadas, la frecuencia de uso, las finalidades académicas, la utilidad percibida y las dificultades reportadas.

Las opciones de respuesta deberán redactarse de forma neutral y comprensible. Cuando corresponda, las categorías serán mutuamente excluyentes; en preguntas de selección múltiple se permitirá marcar varias opciones. También se podrán incluir las alternativas \textit{Otra} y \textit{Prefiero no responder} cuando resulten pertinentes.

\subsection{Instrumento: cuestionario}

El cuestionario se organizará en cinco secciones:

\begin{enumerate}
    \item \textbf{Elegibilidad y consentimiento:} verificación de los criterios de inclusión y aceptación voluntaria de la participación.
    
    \item \textbf{Características académicas:} nivel de formación, modalidad de estudio y semestre o periodo cursado, cuando estos datos sean necesarios para los objetivos.
    
    \item \textbf{Herramientas y frecuencia:} identificación de las categorías de herramientas utilizadas y periodicidad de uso.
    
    \item \textbf{Finalidades académicas:} actividades para las que se emplean las herramientas, como consultar fuentes, redactar documentos, comunicarse, colaborar y organizar tareas.
    
    \item \textbf{Utilidad y dificultades:} valoración de la utilidad percibida y registro de barreras relacionadas con conectividad, dispositivos, costos, habilidades, accesibilidad o privacidad.
\end{enumerate}

Se evitarán preguntas que soliciten contraseñas, credenciales de acceso o información personal innecesaria. Antes de la aplicación definitiva, se realizará una prueba piloto para evaluar la comprensión de los ítems, la pertinencia de las opciones de respuesta y el tiempo requerido para completar el cuestionario. Los ajustes se documentarán antes de iniciar la recolección principal.

\subsection{Técnica complementaria: revisión documental}

Se utilizará la revisión documental para contextualizar el fenómeno y fundamentar el marco referencial. Se consultarán fuentes académicas, informes institucionales y documentos normativos pertinentes para el objeto de estudio.

Además, permitirá orientar la comunicación
de los resultados del estudio. La organización clara y concisa de la información científica es relevante para presentar los hallazgos de manera comprensible y estructurada
\parencite{osmarueda2012resumen}.

Esta actividad no se presentará como una revisión sistemática, dado que la propuesta no contempla un protocolo exhaustivo de búsqueda, criterios de elegibilidad registrados ni una evaluación formal del riesgo de sesgo. Las entrevistas no son indispensables para el objetivo descriptivo actual; podrían incorporarse en una fase posterior si se decide profundizar en las experiencias de los estudiantes.



\subsection{Plan de análisis de datos}

El análisis se desarrollará mediante las siguientes etapas:

\begin{enumerate}
    \item \textbf{Depuración:} revisar la elegibilidad de los registros, identificar duplicados detectables y examinar los datos faltantes. No se imputarán automáticamente respuestas ausentes; se documentará el tratamiento aplicado.

    \item \textbf{Codificación:} asignar códigos a las categorías de respuesta y establecer reglas para el procesamiento de preguntas de selección múltiple. Se conservará un registro de las decisiones de codificación para facilitar la trazabilidad del análisis.

    \item \textbf{Estadística descriptiva:} calcular frecuencias absolutas y porcentajes para las herramientas utilizadas, las finalidades académicas y las dificultades reportadas. La frecuencia de uso y la utilidad percibida se resumirán mediante distribuciones por categoría. Las medidas de resumen se elegirán de acuerdo con el nivel de medición de cada variable.

    \item \textbf{Comparación exploratoria:} cruzar variables seleccionadas, por ejemplo, modalidad de estudio y frecuencia de uso, únicamente cuando el número de observaciones por categoría permita una interpretación prudente. Estas comparaciones no se interpretarán como evidencia de causalidad.

    \item \textbf{Presentación de resultados:} elaborar tablas y gráficos con títulos descriptivos, denominadores explícitos y el número de respuestas válidas correspondiente a cada análisis. Se informarán las limitaciones de los datos y las decisiones de procesamiento.
\end{enumerate}

\subsection{Criterios de interpretación}

Los resultados se interpretarán de acuerdo con el procedimiento de selección de la muestra, la calidad de las respuestas y los objetivos descriptivos. La representatividad no se establecerá únicamente a partir del número de participantes, sino considerando el procedimiento de selección y las características del marco muestral \parencite{argibay2009muestra}.

No se anticiparán valores de significación estadística, intervalos de confianza ni asociaciones como si ya se hubieran calculado. Si se selecciona una muestra probabilística y los objetivos lo justifican, podrán definirse procedimientos inferenciales antes del análisis, con la asesoría metodológica correspondiente. En cambio, si se utiliza una muestra por conveniencia, las comparaciones se interpretarán con prudencia y no se generalizarán automáticamente a toda la población universitaria de Bogotá.

\subsection{Consideraciones éticas y protección de datos}

La participación será voluntaria y los estudiantes podrán retirarse sin penalización. Antes de responder, recibirán información sobre la finalidad del estudio, el uso previsto de los datos, su conservación y las personas autorizadas para consultarlos. Se recopilará únicamente la información necesaria y se limitará el acceso a los archivos de trabajo.

Los resultados se presentarán de forma agregada, evitando la divulgación de datos que permitan identificar a los participantes. La Ley 1581 de 2012 establece disposiciones generales para la protección de datos personales en Colombia \parencite{colombia2012ley1581}; por ello, antes de iniciar la recolección se revisarán las obligaciones aplicables y los procedimientos éticos y de protección de datos de la institución participante.

Si se modifica la población para incluir menores de edad, será necesario revisar el protocolo y obtener las autorizaciones y salvaguardas correspondientes antes de su participación.
